\section{Results}
\label{sec:results}

\subsection{The compiled estate}

The TSA~24 dataset compiles to \num{458708} blocks and 48 treatments; the policy dataset adds the
three prescriptions. The compiled managed area matches the published timber harvesting land base to
\SI{0.32}{\percent} (\autoref{tab:meth_studyregion_spec}).

\subsection{The sustained harvest level}

The even-flow search converges to realised levels of \SIrange{51.89}{52.06}{\mega\cubic\meter}
per period, a mean of \SI{5.19}{\mega\cubic\meter\per\year}; the replayed schedule ranges from
\SI{51.2}{\mega\cubic\meter} to \SI{52.3}{\mega\cubic\meter} across periods. We treat this as
the estate's sustained even-flow level, not a proven maximum.

The level is about \SI{25}{\percent} below the official allowable annual cut of
\SI{6917231}{\cubic\meter\per\year}. The managed area matches the published land base to within
\SI{0.32}{\percent} (\SI{3060473}{\hectare} against \SI{3070301}{\hectare}), and the model's
own curves imply \SI{7.07}{\mega\cubic\meter\per\year} at their area-weighted maximum mean
annual increment, about \SI{2}{\percent} above the cut. The shortfall is therefore
consistent with the flat-flow rule on a non-regulated age structure, not with the land base or the
growth potential of the curves. The official determination is a declining transition profile; a
constant profile cannot take the transition harvest that a declining profile allows.

\subsection{Economic value of the model-produced schedule}

Valued with the constants in \autoref{sec:methods}, the model-produced schedule harvests
\SI{1037.9}{\mega\cubic\meter} over the horizon, produces \SI{112776}{million\CAD} of log revenue
and \SI{62270}{million\CAD} of undiscounted net value. At a \SI{3}{\percent} real discount rate the
net present value is \SI{8779}{million\CAD} (\autoref{tab:res_economics}).

Unit margins differ sharply by species group. Conifers return between
\SI{56.81}{CAD\per\cubic\meter} and \SI{64.79}{CAD\per\cubic\meter} after harvest, haul,
silviculture, and stumpage. Deciduous species return \SI{-2.25}{CAD\per\cubic\meter}: at the
recorded interior hardwood price, harvesting deciduous timber does not cover its delivered cost.
Deciduous species are \SI{5.4}{\percent} of harvested volume in this schedule, so the schedule
carries a small negative-margin component, which any composition policy that raises the deciduous
share will enlarge.

\begin{table}[h]
    \centering
    \caption{Economic result of the model-produced schedule over 20 periods of ten years.}
    \label{tab:res_economics}
    \small
    \begin{tabular}{lr}
        \toprule
        Item & Value \\
        \midrule
        Harvested volume                   & \SI{1037.9}{\mega\cubic\meter} \\
        Log revenue                        & \SI{112776}{million\CAD} \\
        Undiscounted net value             & \SI{62270}{million\CAD} \\
        Net present value, \SI{3}{\percent}& \SI{8779}{million\CAD} \\
        Conifer unit margin                & \SIrange{56.81}{64.79}{CAD\per\cubic\meter} \\
        Deciduous unit margin              & \SI{-2.25}{CAD\per\cubic\meter} \\
        \bottomrule
    \end{tabular}
\end{table}

\subsection{Sensitivity to the documented ranges}

Three economic constants are ranges rather than points: the discount rate, the deciduous log
price, and the haul distance. Their effect on the model-produced schedule is shown in
\autoref{tab:res_sensitivity}.

The discount rate dominates: moving from \SI{3}{\percent} to \SI{5}{\percent} reduces net
present value by \SI{4010}{million\CAD}, close to half the base value. No other constant moves the
result by more than about three percent. The deciduous price moves net present value by
\SI{26}{million\CAD} in either direction; the haul distance from \SI{90}{\kilo\meter} to
\SI{117}{\kilo\meter} costs \SI{184}{million\CAD}.

\begin{table}[h]
    \centering
    \caption{Sensitivity of net present value (base \SI{8779}{million\CAD}) to the documented
    ranges of the economic constants.}
    \label{tab:res_sensitivity}
    \small
    \begin{tabular}{llrr}
        \toprule
        Lever & Setting & NPV & Change \\
        \midrule
        Discount rate & \SI{5}{\percent} & \SI{4769}{million\CAD} & \SI{-4010}{million\CAD} \\
        Deciduous price & \SI{35}{CAD\per\cubic\meter} & \SI{8752}{million\CAD} & \SI{-26}{million\CAD} \\
        Deciduous price & \SI{40}{CAD\per\cubic\meter} & \SI{8805}{million\CAD} & \SI{+26}{million\CAD} \\
        Haul distance & \SI{117}{\kilo\meter} & \SI{8595}{million\CAD} & \SI{-184}{million\CAD} \\
        \bottomrule
    \end{tabular}
\end{table}

Two readings follow. First, the economic result is robust to the constants that are genuinely
uncertain, and sensitive to the one that is a policy choice. Second, the mean value is not where a
composition policy will show its effect: the deciduous price range, which spans the plausible
value of the species that a diversification policy would promote, moves the result by a fraction
of a percent. If composition matters, it matters in the tail.

\subsection{Stocking standards and the no-fire surface}

The historical clearcut regenerates the same analysis unit. On a conifer site it returns a
conifer stand; the aspen share of the future forest is therefore limited by the aspen already on
the land base. The three prescriptions regenerate a target analysis unit with 26, 55, or
84 percent deciduous stock. The instrument changes what is planted, so composition is no longer
limited to the existing aspen supply.

We evaluate the three harvest levels against the four standards with the no-fire process model
(\autoref{tab:res_standards}). The CC row is the no-fire baseline (\SI{8778.8}{\mega\CAD},
\SI{1037.9}{\mega\cubic\meter}); the other rows are compared with it.

\begin{table}[h]
    \centering
    \small
    \caption{No-fire production and value by harvest level and stocking standard, 200 years. The
    delta columns are relative to CC at the same harvest level. NPV is in million CAD and volume
    in million cubic metres.}
    \label{tab:res_standards}
    \begin{tabular}{lrrrrr}
        \toprule
        Level & Standard & NPV & Volume & $\Delta$NPV & $\Delta$Volume \\
        \midrule
        100\% & CC  & 8,778.8 & 1,037.9 & 0.0 & 0.0 \\
        100\% & D26 & 8,779.8 & 1,091.2 & +1.0 & +53.3 \\
        100\% & D55 & 8,671.2 & 1,016.0 & -107.6 & -21.9 \\
        100\% & D84 & 8,625.0 & 1,014.1 & -153.8 & -23.8 \\
        \midrule
        80\%  & CC  & 7,015.6 & 829.7 & 0.0 & 0.0 \\
        80\%  & D26 & 7,016.5 & 872.2 & +0.8 & +42.5 \\
        80\%  & D55 & 6,929.6 & 812.2 & -86.0 & -17.5 \\
        80\%  & D84 & 6,892.6 & 810.7 & -123.0 & -19.0 \\
        \midrule
        60\%  & CC  & 5,247.6 & 620.7 & 0.0 & 0.0 \\
        60\%  & D26 & 5,248.1 & 652.6 & +0.5 & +32.0 \\
        60\%  & D55 & 5,182.4 & 607.4 & -65.2 & -13.2 \\
        60\%  & D84 & 5,154.2 & 606.3 & -93.4 & -14.4 \\
        \bottomrule
    \end{tabular}
\end{table}

The D26 standard is \emph{free volume}: it adds \SIrange{32}{53}{\mega\cubic\meter} and does not
reduce NPV. The aspen yield curves add volume, and the low aspen price offsets the value. Under
the calibrated 30 percent mixedwood threshold the D26 stand stays on the conifer side (C2/C1), so
it buys no fire benefit. It is a production option, not a risk option.

The D55 standard costs about \SI{1.2}{\percent} of NPV and loses about \SI{2.1}{\percent} of
volume. Its fuel type is M1 and its severity multiplier is 1.011, so it is also a production
option, dominated by D26. The D84 standard costs about \SI{1.8}{\percent} of NPV and loses about
\SI{2}{\percent} of volume. Its fuel type is D1 and its severity multiplier is 0.267. It is the
only standard with a material fire benefit.

\subsection{Fire regime}

The policy surface uses the full recorded regime: \num{1073} fires from 1950 to 2024,
\num{14.31} fires per year, and \SI{10079}{\hectare\per\year}. The calibrated spread model is
compared with the 432-perimeter fitting sample in \autoref{tab:res_fire}; the median and the share
above \SI{500}{\hectare} sit inside the historical interval, the 90th percentile at its lower
edge, and the 95th percentile and the extreme-tail shares below it. The central distribution is
reproduced, and the upper tail is thin relative to the historical record. The maximum is reported
for completeness but not used as a target: its bootstrap interval is one-sided.

\begin{table}[h]
    \centering
    \small
    \caption{Simulated fire-size distribution against the historical record. Historical intervals
    are the 5th to 95th percentiles of 20,000 bootstrap resamples of the 432 perimeters; simulated
    values are point estimates from 3,450 fires and were not bootstrapped. The maximum interval is
    one-sided and is not used as a target.}
    \label{tab:res_fire}
    \begin{tabular}{lrrr}
        \toprule
        Statistic & Historical & Historical interval & Simulated \\
        \midrule
        Median [ha]            & 22.97   & 20.85 to 30.16      & 23.8 \\
        90th percentile [ha]   & 490.17  & 349.13 to 1,051.26  & 361.5 \\
        95th percentile [ha]   & 1,918.88& 1,415.35 to 2,955.33& 1,100.9 \\
        Share above 100 [\%]   & 26.2    & 22.7 to 29.6        & 23.6 \\
        Share above 500 [\%]   & 10.0    & 7.6 to 12.5         & 8.3 \\
        Share above 1,000 [\%] & 7.9     & 5.8 to 10.2         & 5.5 \\
        Share above 5,000 [\%] & 3.0     & 1.9 to 4.4          & 1.8 \\
        Maximum [ha]           & 25,114.8& one-sided           & 20,770.4 \\
        \bottomrule
    \end{tabular}
\end{table}

The calibrated model is also run on the full estate with 24 ignitions to confirm that it couples to
the harvest model: it grows \SI{16331}{\hectare} and applies \SI{10444}{\hectare}, with an ending
volume of \SI{1015.9}{\mega\cubic\meter}. The run is a coupling check, not a distribution
sample.

\subsection{The fire-risk surface}

The paired ensemble gives the risk side (\autoref{tab:res_fire_surface}). The CC row is the
baseline: \SI{243.8}{\mega\CAD} mean loss and \SI{304.3}{\mega\CAD} CVaR20.

\begin{table}[h]
    \centering
    \small
    \caption{Fire loss by harvest level and stocking standard. Ten paired draws of the full
    \num{1073}-fire regime, 200 years, composition-aware severity. CVaR20 is the mean of the
    two worst draws; the percentage is relative to the no-fire net present value.}
    \label{tab:res_fire_surface}
    \begin{tabular}{lrrrr}
        \toprule
        Level & Standard & Mean loss [M CAD] & CVaR20 [M CAD] & CVaR20 [\%] \\
        \midrule
        100\% & CC  & 243.8 & 304.3 & 3.47 \\
        100\% & D26 & 243.8 & 304.3 & 3.47 \\
        100\% & D55 & 238.5 & 298.6 & 3.44 \\
        100\% & D84 & 236.3 & 296.2 & 3.43 \\
        \midrule
        80\%  & CC  & 193.2 & 239.7 & 3.42 \\
        80\%  & D26 & 193.3 & 239.7 & 3.42 \\
        80\%  & D55 & 189.1 & 235.2 & 3.39 \\
        80\%  & D84 & 187.3 & 233.2 & 3.38 \\
        \midrule
        60\%  & CC  & 144.7 & 179.3 & 3.42 \\
        60\%  & D26 & 144.7 & 179.4 & 3.42 \\
        60\%  & D55 & 141.5 & 175.9 & 3.39 \\
        60\%  & D84 & 140.2 & 174.4 & 3.38 \\
        \bottomrule
    \end{tabular}
\end{table}

The D26 standard is not a risk instrument. Its stand composition stays on the conifer side of
the calibrated 30 percent mixedwood threshold, and its CVaR20 is indistinguishable from CC
(+0.01 M at the 100 percent level). It is a free-volume option with a substantial composition change:
the landscape deciduous share rises from 7.5 to 16.7 percent without changing fire severity.

The D55 and D84 standards are a priced tail hedge. D55 cuts CVaR20 by \SI{5.7}{\mega\CAD} at a
production cost of \SI{107.6}{\mega\CAD}; D84 cuts it by \SI{8.1}{\mega\CAD} at
\SI{153.8}{\mega\CAD}. The paired 95 percent bootstrap intervals are \SIrange{-6.2}{-4.1}{\mega\CAD}
for D55 and \SIrange{-8.8}{-5.9}{\mega\CAD} for D84 at the 100 percent level; both exclude zero.
The exchange rate is about \textbf{19 M CAD of production cost per 1 M CAD of CVaR20 reduction}, and it is stable across levels (18.9 to 19.2) and standards. The
D84 stand-level severity multiplier is 0.267, but the estate mean loss falls only from 243.8 to
236.3 M CAD (3.1 percent) and CVaR20 from 304.3 to 296.2 M CAD (2.7 percent). The reasons are
structural: the unmanaged half of the estate keeps its VRI fuel, the standard converts stands
gradually over the horizon, the conversion begins with a 25-year slash/open phase, and the NPV
loss is dominated by early, discounted fires.

A risk-neutral planner therefore chooses D26: more volume, the same NPV, and the same fire
risk. A planner who values a \SI{1}{\mega\CAD} reduction in CVaR20 at more than about
\SI{19}{\mega\CAD} chooses D84. D55 lies on the same frontier as D84 and is an intermediate
point, not a distinct strategy. A more severe future regime, a larger managed fraction, a higher
deciduous price, or a carbon value would move the exchange rate.

\subsection{Rotation length and diversification}

We estimate the rotation from the schedule: for every block harvested more than once, the interval
between successive harvests. The baseline yields \num{53071} intervals, an area-weighted mean
rotation of \SI{102.2}{\year}, with a median of \SI{100}{\year} and a tenth-to-ninetieth
percentile range of \SIrange{60}{160}{\year}; \SI{44}{\percent} of the managed area is older
than 100 years. The estimate is conditional on re-entry within 200 years and is biased toward
productive and accessible stands. The new standards are evaluated on a fixed schedule, so their
rotation is an outcome, not a re-optimised choice.

The rotation is consistent with the level: at 102 years over \SI{3.06}{\mega\hectare}, about
\SI{30000}{\hectare} are harvested each year; at \SI{5.19}{\mega\cubic\meter\per\year} that
is \SI{173}{\cubic\meter\per\hectare}, about \SI{73}{\percent} of the
\SI{236}{\cubic\meter\per\hectare} that the area-weighted maximum mean annual increment
accumulates over the rotation.
